\section*{Chapitre \ref{ch:iv3:reading-debugging-and-software-design} --- Lecture, débogage et conception logicielle}
\markboth{\protect\omTocMark{}{Réponses~: Lecture, débogage et conception logicielle}}{\protect\omTocMark{}{Réponses~: Lecture, débogage et conception logicielle}}

\begin{solution}{iq:iv3:reading-debugging-and-software-design:1}
Le programme affiche \texttt{more bids than asks}, puis \texttt{18446744073709551615}. Rien n'est indéfini. La correction consiste à comparer directement les tailles~: \texttt{bids.size() > asks.size()}.

\texttt{size()} renvoie \texttt{std::size\_t}, un type non signé de 64 bits ici. La norme définit l'arithmétique non signée modulo $2^{64}$, donc $2 - 3$ vaut $2^{64} - 1$, qui est la deuxième ligne de sortie. Une valeur non signée n'est jamais négative, donc \texttt{> 0} est vrai pour toute paire de tailles différentes, quelle que soit la plus grande. Sous \texttt{-Wall -Wextra -Werror}, le listing compile sans le moindre avertissement, car \texttt{> 0} sur une valeur non signée est un test raisonnable en général.

Si l'on a vraiment besoin de la différence signée, C++20 fournit \texttt{std::ssize}. Soustraire deux de ces valeurs et tester \texttt{> 0} donne aussi la bonne branche.

\iqwrong{«~Le dépassement de capacité est un comportement indéfini, donc tout peut arriver.~» C'est vrai pour les types signés. L'arithmétique non signée reboucle par définition, donc ce programme est parfaitement défini et parfaitement faux.}

\iqfollow{Que fait \texttt{for (std::size\_t i = v.size() - 1; i >= 0; -{}-i)}~?}{La boucle ne se termine jamais, car \texttt{i >= 0} est toujours vrai~: après 0, elle reboucle sur la plus grande valeur et indexe bien au-delà des bornes. Ici \texttt{-Wextra} la signale pourtant, sous la forme de \texttt{-Wtype-limits}. Avec un \texttt{v} vide, le programme échoue encore plus tôt, dès la première itération.}

\iqwhy{Le calcul sur les quantités et les niveaux de prix est plein de tailles. Ce bogue passe tous les tests où le premier conteneur se trouve être le plus grand.}
\end{solution}

\begin{solution}{iq:iv3:reading-debugging-and-software-design:2}
Il affiche \texttt{[30, 30, 30]}. Pour obtenir \texttt{[10, 20, 30]}, il faut figer la valeur au moment de la création de la fonction~: \texttt{lambda x, k=k: k * x}. Une compréhension de liste n'y change rien.

Une fermeture en Python capture la variable, et non sa valeur. Les trois lambdas renvoient à une même variable \texttt{k} dans la portée de \texttt{make\_scalers}, et elles la consultent au moment de l'appel. À ce moment-là, la boucle est terminée et \texttt{k} vaut 3. Les arguments par défaut sont évalués à la définition de la fonction, ce qui fait que \texttt{k=k} prend un instantané. \texttt{functools.partial} fait la même chose et ne laisse aucun paramètre supplémentaire qu'un appelant pourrait fournir par erreur.

\iqwrong{«~Il faut utiliser une compréhension, puisqu'elle a sa propre portée.~» Elle en a une, mais une seule pour toute la compréhension. \texttt{[lambda x: k * x for k in [1, 2, 3]]} donne aussi \texttt{[30, 30, 30]}~: le test l'exécute.}

\iqwhy{Construire une liste de callbacks ou de scalers dans une boucle est du code de recherche courant, et l'échec est silencieux~: chaque instrument reçoit les paramètres du dernier instrument.}
\end{solution}

\begin{solution}{iq:iv3:reading-debugging-and-software-design:3}
Sur ce compilateur, il affiche \texttt{10 0 20}~: un zéro subsiste, et le programme a un comportement indéfini. Si un zéro se trouve en dernière position, la boucle sort de la fin du vecteur. Il faut utiliser \texttt{std::erase(qty, 0L)}.

\texttt{erase} invalide tous les itérateurs à partir de l'élément effacé, si bien que \texttt{it} ne peut plus servir, et \texttt{++it} est déjà un comportement indéfini. À l'exécution, les éléments ont glissé vers la gauche, donc \texttt{it} pointait sur le second zéro, et \texttt{++it} l'a sauté. C'est pour cela qu'un zéro subsiste.

Avec l'entrée \texttt{\{10, 0\}}, l'effacement laisse \texttt{it} égal à \texttt{end()}, le \texttt{++it} le dépasse, et le test \texttt{it != qty.end()} ne devient jamais faux. AddressSanitizer signale un heap-buffer-overflow sur cette entrée. Le mode de débogage de libstdc++, \texttt{-D\_GLIBCXX\_DEBUG}\index{mode de débogage, libstdc++}, arrête même la première entrée, avec un message sur un itérateur singulier.

\texttt{std::erase} de C++20 fait le travail en un seul passage. Avant C++20, on écrirait l'idiome erase-remove. La boucle manuelle peut être corrigée avec \texttt{it = qty.erase(it)} et sans incrément dans cette branche, mais chaque effacement décale les éléments suivants, si bien qu'elle coûte $O(n^2)$ là où \texttt{std::erase} coûte $O(n)$.

\iqwrong{«~Il retire tous les zéros, puisque \texttt{erase} retourne normalement.~» En lisant la boucle sur \texttt{\{10, 0, 0, 20\}} pas à pas, on voit tout de suite le saut. Il faut le dire à voix haute pendant l'entretien.}

\iqwhy{Supprimer des éléments pendant un parcours est la question de lecture classique, et la relance sur le dernier élément permet de voir si le candidat repère le comportement indéfini ou seulement la sortie fausse.}
\end{solution}

\begin{solution}{iq:iv3:reading-debugging-and-software-design:4}
Il affiche \texttt{1}, \texttt{1}, \texttt{2}. La mauvaise réponse~: un symbole sans limite est autorisé. L'effet de bord~: la consultation de \texttt{BBB} l'insère, si bien que le tableau associatif passe d'une entrée à deux.

\texttt{std::map::operator[]} est défini comme \texttt{try\_emplace(x).first->second}~: si la clé est absente, il insère une entrée initialisée par valeur, ici 0, et renvoie une référence sur celle-ci. Ainsi \texttt{limits[\textquotedbl{}BBB\textquotedbl{}]} crée une limite de 0, et $0 \le 0$ passe. Un contrôle de risque qui répond oui à un instrument dont il n'a jamais entendu parler doit échouer en position sûre.

Pour que le compilateur le détecte, il suffit de rendre le tableau associatif \texttt{const}, ou de le passer par \texttt{const\&}. \texttt{operator[]} n'a pas de surcharge const, donc le code ne compile plus~: c'est ce que montre le test. Il faut ensuite écrire la recherche avec \texttt{find}~:

\texttt{const auto it = limits.find(symbol);}

\texttt{return it != limits.end() \&\& qty <= it->second;}

Cela affiche \texttt{1}, \texttt{0}, \texttt{1}. \texttt{at} est l'autre choix, et il lève une exception sur une clé absente.

\iqfollow{Pourquoi l'effet de bord compte-t-il si la mauvaise réponse est corrigée autrement~?}{Parce que le tableau associatif est écrit pendant ce qui ressemble à une lecture. Deux threads qui appellent \texttt{allowed} créent maintenant une course aux données, alors qu'aucun des deux ne se croit écrivain. La mémoire grossit aussi avec chaque symbole mal orthographié.}

\iqwhy{C'est l'exemple le plus court d'une lecture qui écrit, et une ligne de ce genre figure dans de vrais contrôles de risque.}
\end{solution}

\begin{solution}{iq:iv3:reading-debugging-and-software-design:5}
Il affiche \texttt{start base}, puis \texttt{run momentum}. Pendant le constructeur de \texttt{Strategy}, l'objet est un \texttt{Strategy}, donc l'appel va vers \texttt{Strategy::name}.

La norme dit qu'une fonction virtuelle appelée depuis un constructeur se résout vers la redéfinition finale dans la classe du constructeur, et non dans une classe plus dérivée. La raison est valable~: les membres de \texttt{Momentum} n'existent pas encore, et une redéfinition pourrait s'en servir. Une fois la construction terminée, \texttt{m.name()} voit \texttt{Momentum} comme prévu.

Si \texttt{name} est une fonction virtuelle pure et que le constructeur l'appelle directement, le compilateur peut le signaler. Passé par une fonction auxiliaire, il ne voit pas le problème. Cet appel a un comportement indéfini, et avec \texttt{-O0} l'exécution du test s'interrompt avec \texttt{pure virtual method called}.

Pour obtenir ce que voulait l'auteur, il faut séparer la construction du démarrage. Soit prévoir une fabrique qui construit l'objet puis appelle \texttt{start()} dessus, soit faire passer le nom par le constructeur~: \texttt{Strategy(const char* name)}.

\iqwrong{«~Les appels virtuels vont toujours à la classe la plus dérivée.~» Ce n'est pas vrai pendant la construction ni la destruction. Il en va de même dans un destructeur, où la partie dérivée a déjà disparu.}

\iqwhy{Journaliser ou enregistrer quelque chose dans le constructeur d'une classe de base est une chose naturelle à écrire, et c'est la façon habituelle dont ce bogue apparaît~: une stratégie s'enregistre sous le mauvais nom.}
\end{solution}

\begin{solution}{iq:iv3:reading-debugging-and-software-design:6}
Il affiche \texttt{0.0 101.0 0.0}. C'est une troncature d'objet\index{troncature d'objet}. Il faut stocker des pointeurs dans le vecteur, en général des \texttt{unique\_ptr}, et interdire le constructeur de copie de la classe de base pour que l'erreur soit rejetée dès la compilation.

Un vecteur d'\texttt{Instrument} contient des objets de ce type exact, chacun avec la place d'un seul \texttt{Instrument}. \texttt{push\_back(Bond(99.5))} ne copie que la partie \texttt{Instrument}, si bien que le prix disparaît et que le pointeur de table virtuelle est celui de \texttt{Instrument}~: \texttt{book[0].value()} vaut 0.0. \texttt{copy} est la même opération écrite explicitement~: \texttt{Instrument copy = b}. Seul \texttt{ref} est encore un \texttt{Bond}, donc il donne 101.0.

Avec le vecteur de \texttt{unique\_ptr}, le premier nombre devient 99.5. Il suffit d'ajouter \texttt{Instrument(const Instrument\&) = delete;}, avec un constructeur par défaut demandé explicitement, et les deux lignes de troncature cessent de compiler. C'est la règle habituelle pour une classe destinée à être utilisée par un pointeur de base~: pas de copie publique. Si l'on a vraiment besoin de copies, il faut ajouter une fonction virtuelle \texttt{clone}.

\iqfollow{Existe-t-il un moyen de garder les instruments côte à côte en mémoire~?}{Oui, si l'ensemble des types est fermé~: \texttt{std::vector<std::variant<Bond, Swap>{}>} stocke chaque valeur sur place, sans troncature ni allocation supplémentaire. Le coût est qu'ajouter un type oblige à modifier cette ligne.}

\iqwhy{Le code de valorisation d'une banque regorge de hiérarchies d'instruments, et une troncature donne un nombre plausible sans la moindre erreur.}
\end{solution}

\begin{solution}{iq:iv3:reading-debugging-and-software-design:7}
Sous \texttt{cargo run}, il panique, et le message est \texttt{attempt to subtract with overflow}. Avec \texttt{-{}-release}, il affiche $2^{64} - 1$ en entier. Aucun des deux cas n'est unsafe au sens de Rust. Il faut utiliser \texttt{abs\_diff} pour la taille de l'écart, ou un résultat signé si le signe compte.

\texttt{len()} renvoie \texttt{usize}, et $2 - 3$ ne tient pas. The Rust Reference qualifie le dépassement arithmétique d'une erreur de programmeur, mais ce n'est pas un comportement indéfini. Avec les assertions de débogage activées, le compilateur doit insérer une vérification qui panique. En mode release, le comportement par défaut est de reboucler, comme le confirme le test. Le programme reste sûr en mémoire dans les deux cas~: soit le mauvais nombre sort, soit le processus s'arrête, et aucun invariant du langage n'est rompu. On peut conserver les vérifications en mode release avec \texttt{overflow-checks = true} dans le profil Cargo, et le test confirme qu'une construction optimisée panique alors aussi.

La correction dépend de l'appelant. \texttt{bids.len().abs\_diff(asks.len())} donne 1 ici si seule compte la taille de l'écart. \texttt{checked\_sub} renvoie \texttt{None} si un résultat négatif signifie «~cela ne devrait pas arriver~». \texttt{a as i64 - b as i64} donne $-1$ si l'appelant veut le signe.

\iqwrong{«~Rust détecte les dépassements de capacité.~» Seulement avec les vérifications activées. Le même code livre un comportement différent de celui qu'on a testé, donc le test en mode release compte.}

\iqwhy{Les sociétés crypto et les teneurs de marché qui utilisent Rust posent cette question pour vérifier qu'on sait ce que safe garantit et ce qu'il ne garantit pas.}
\end{solution}

\begin{solution}{iq:iv3:reading-debugging-and-software-design:8}
Il affiche \texttt{False 0}, puis \texttt{2.67 0 2 2}. Il faut dire à son collègue d'indexer les niveaux de prix par un nombre entier de pas de cotation.

$0.1 + 0.2$ vaut \texttt{0.30000000000000004} en virgule flottante binaire, et le littéral \texttt{0.3} est un autre double. La recherche dans un dict utilise l'égalité exacte, donc l'entrée existe sous une clé et la recherche en demande une autre, et l'exécution semble disparaître. Deux prix qui s'affichent de la même façon peuvent différer du dernier bit s'ils proviennent de calculs différents.

La deuxième ligne montre deux autres pièges dans \texttt{round}. \texttt{2.675} est stocké un peu en dessous de cette valeur, sous la forme \texttt{2.67499999\dots}, donc il s'arrondit vers le bas, ce que la documentation Python signale elle-même. Les cas de moitié vont au voisin pair~: \texttt{round(0.5)} vaut 0, \texttt{round(1.5)} et \texttt{round(2.5)} valent tous deux 2. Arrondir au nombre de décimales du pas de cotation donne bien un double canonique, donc \texttt{round(0.1 + 0.2, 2) == 0.3}, mais cela ne marche que tant que l'erreur reste bien en deçà d'un demi-pas de cotation, et chaque auteur de code doit y penser. Des pas entiers n'exigent pas une telle précaution.

Les pas entiers le résolvent complètement~: avec un pas de cotation de 0.001, les clés 100 et 200 s'additionnent exactement à 300. \texttt{Decimal(\textquotedbl{}0.1\textquotedbl{}) + Decimal(\textquotedbl{}0.2\textquotedbl{})} vaut aussi \texttt{Decimal(\textquotedbl{}0.3\textquotedbl{})}, ce qui convient aux rapports et aux montants monétaires, au prix d'une vitesse moindre.

\iqwrong{«~Il suffit de comparer avec une tolérance.~» On ne peut pas chercher dans un dict avec une tolérance, et la relation de tolérance n'est pas transitive, donc on ne peut pas non plus en tirer une clé cohérente.}

\iqwhy{Représenter les prix par des flottants est le défaut de conception le plus courant dans les carnets d'ordres des candidats, et de nombreuses sociétés le vérifient dans la première minute.}
\end{solution}

\begin{solution}{iq:iv3:reading-debugging-and-software-design:9}
Il affiche \texttt{true false true} puis \texttt{true false true}. Le langage garantit chacun de ces résultats sauf \texttt{c == d}~: seul ce false est propre à cette JVM.

\texttt{==} sur des références demande s'il s'agit du même objet. La Java Language Specification précise que la boxe d'un entier compris entre $-128$ et 127, effectuée deux fois, donne le même objet, si bien que \texttt{a == b} est toujours vrai. Au-delà de cet intervalle, rien n'est garanti dans un sens comme dans l'autre. La spécification ajoute que les implémentations moins contraintes en mémoire peuvent mettre en cache davantage, jusqu'à $\pm$32K pour \texttt{int}. OpenJDK 11 avec les réglages par défaut donne deux objets pour 128, donc \texttt{c == d} est faux.

Pour les chaînes, la spécification dit que les littéraux sont internés, donc \texttt{s == t} est vrai. \texttt{new String(\textquotedbl{}AAA\textquotedbl{})} crée par définition un nouvel objet, donc \texttt{s == u} est faux, alors que \texttt{equals} compare les contenus.

La règle tient en peu de mots~: utiliser \texttt{equals} sur les objets, et \texttt{==} seulement sur les primitifs. Ce bogue est vicieux, car il passe tous les tests avec de petites quantités.

\iqwrong{«~\texttt{==} sur \texttt{Integer} compare les valeurs.~» Ce serait le cas si l'un des deux côtés était un \texttt{int} primitif, car l'autre côté est alors déballé. Avec deux variables \texttt{Integer}, on compare des références.}

\iqwhy{Les systèmes bancaires regorgent de valeurs boxées dans des collections. Une vérification du type \texttt{qty == limit} passe dans tous les petits cas de test et casse à 128.}
\end{solution}

\begin{solution}{iq:iv3:reading-debugging-and-software-design:10}
Il affiche \texttt{100}, puis \texttt{\textasciitilde{}Order}, et \texttt{\textasciitilde{}LimitOrder} ne s'exécute jamais. Le programme a un comportement indéfini~: supprimer un \texttt{LimitOrder} par un \texttt{Order*} exige un destructeur virtuel. AddressSanitizer l'arrête. Avec \texttt{std::make\_shared<LimitOrder>}, il est défini, et les deux destructeurs s'exécutent.

La norme dit que, lors de la suppression d'un objet unique, si le type statique diffère du type dynamique, le type statique doit posséder un destructeur virtuel, faute de quoi le comportement est indéfini. \texttt{qty} est une fonction virtuelle, donc la classe est polymorphe, mais le destructeur ne l'est pas. Dans cette exécution, \texttt{tag} n'est jamais détruit et son tampon alloué sur le tas est perdu.

ASan signale \texttt{new-delete-type-mismatch}~: le type alloué fait 40 octets, et le type libéré 8. Avec ce contrôle désactivé, il signale plutôt la fuite de la chaîne. Le test montre les deux.

Pourquoi aucun avertissement n'a-t-il été émis~? g++ possède \texttt{-Wdelete-non-virtual-dtor}, que \texttt{-Wall} active, mais l'expression \texttt{delete} se trouve ici dans \texttt{unique\_ptr}, dans un en-tête système, où les avertissements sont supprimés. Avec \texttt{delete p} sur un pointeur brut, la compilation échoue avec les mêmes options, comme le confirme le test. Il ne faut donc pas se fier au silence.

La correction tient en un mot~: \texttt{virtual \textasciitilde{}Order()}. Le programme affiche alors \texttt{\textasciitilde{}LimitOrder}, puis \texttt{\textasciitilde{}Order}, et ASan reste silencieux.

\iqfollow{Pourquoi \texttt{std::make\_shared<LimitOrder>} fonctionne-t-il sans destructeur virtuel~?}{Le bloc de contrôle du pointeur partagé mémorise comment détruire l'objet avec lequel il a été créé, à savoir le \texttt{LimitOrder} réel, si bien que le suppresseur ne passe pas par \texttt{Order}. Le test affiche les deux messages de destruction, et ASan reste silencieux. C'est garanti, mais seulement pour ce type de pointeur~: il faut tout de même ajouter le destructeur virtuel, car la personne suivante utilisera \texttt{unique\_ptr}.}

\iqwhy{C'est la question classique des entretiens C++, et le piège tient à ce que le pointeur intelligent habituel masque l'avertissement.}
\end{solution}

\begin{solution}{iq:iv3:reading-debugging-and-software-design:11}
Il fuit un \texttt{Order} à chaque fois que \texttt{validate} lève une exception. Ici, cela représente un objet de 40 octets, et rien d'autre, car le tag court tient dans le tampon interne de la chaîne. LeakSanitizer le détecte. Il faut utiliser \texttt{std::make\_unique}. Ne pas se contenter du \texttt{catch}.

Le chemin en cause est celui de l'ordre \texttt{-1}. \texttt{new} réussit, \texttt{validate} lève une exception, la pile se déroule, et le pointeur brut \texttt{o} est détruit sans toucher à ce qu'il désigne. Donc \texttt{delete} ne s'exécute jamais. Le total reste exact, 300, car la fuite n'a rien à voir avec l'arithmétique, ce qui explique qu'elle ait pu durer des mois~: quelques octets par ordre rejeté.

Le contrôle de fuites d'ASan signale \texttt{Direct leak of 40 byte(s) in 1 object(s)}, alloué dans \texttt{handle}. Avec \texttt{auto o = std::make\_unique<Order>(\dots)} et aucun \texttt{delete}, le destructeur s'exécute pendant le déroulement de la pile. Le test le confirme~: même sortie, rien dans le rapport.

La version du collègue fonctionne bien pour un seul pointeur. Mais elle duplique le nettoyage sur deux chemins. Le problème s'aggrave à chaque nouvelle ressource, car une deuxième allocation exige des gestionnaires imbriqués. Et la personne qui ajoutera plus tard un \texttt{return} anticipé l'oubliera. C'est pourquoi la règle veut que chaque ressource ait un objet propriétaire dont le destructeur la libère. Ici, en fait, aucun tas n'est nécessaire~: \texttt{const Order o\{qty, \textquotedbl{}tag\textquotedbl{}\};} est le code correct le plus simple.

\iqwrong{«~Ce ne sont que 40 octets.~» Il s'agit de 40 octets à chaque fois, et la fuite croît avec le taux de rejets, qui grimpent justement quand le marché ou un client dysfonctionne.}

\iqwhy{Ce que teste la question, c'est l'habitude de lire le chemin qui lève une exception, et non celui qui s'exécute normalement.}
\end{solution}

\begin{solution}{iq:iv3:reading-debugging-and-software-design:12}
Le programme a un comportement indéfini en C++20, et ASan signale \texttt{stack-use-after-scope}. La version avec le membre public est définie et donne 600. Le C++23 rend aussi la première version définie.

Une boucle for sur intervalle est définie par \texttt{auto\&\& range = \textit{initializer};}, suivie d'une boucle. Lier une référence à un temporaire prolonge la durée de vie de ce temporaire, mais seulement si la référence se lie à lui, ou à l'un de ses sous-objets. Ici, elle se lie à ce que renvoie \texttt{levels()}. Cette fonction renvoie une référence, et le compilateur ne voit pas qu'elle pointe dans le temporaire \texttt{Snapshot}. Donc le \texttt{Snapshot} meurt à la fin de l'initialiseur, avant le début de la boucle, et celle-ci parcourt un vecteur détruit.

\texttt{latest().levels\_} désigne directement un sous-objet du temporaire, si bien que la prolongation s'applique à tout le \texttt{Snapshot}. Le test s'exécute proprement sous ASan. Ainsi, cette ligne d'apparence identique est sûre par un membre de données, et dangereuse par un accesseur.

Le C++23 a ajouté une règle à \texttt{[class.temporary]}~: un temporaire créé dans l'initialiseur d'une boucle for sur intervalle vit aussi longtemps que la référence à l'intervalle. g++ 11.4 en mode C++20 ne fait pas cela, comme le montre le rapport d'ASan.

En C++20, il faut garder l'objet dans une variable nommée. Écrire d'abord \texttt{const Snapshot snap = latest();}, puis boucler sur \texttt{snap.levels()}. Pour que le compilateur rejette le code, il faut ajouter des qualificateurs de référence~: \texttt{levels() const\&} accompagné d'une surcharge supprimée \texttt{levels() const\&\&}. Le test montre que le code ne compile alors plus.

\iqwrong{«~Les temporaires vivent jusqu'à la fin de l'instruction, et la boucle est une seule instruction.~» Le temporaire vit jusqu'à la fin de l'expression complète qui initialise l'intervalle, laquelle se termine avant l'exécution du corps de la boucle.}

\iqwhy{Un accesseur appliqué à la valeur de retour d'une fonction, dans une boucle, est une forme que l'on trouve dans toute base de code C++. Elle semble souvent fonctionner, car la mémoire libérée est encore là.}
\end{solution}

\begin{solution}{iq:iv3:reading-debugging-and-software-design:13}
La ligne \texttt{std::lock\_guard<std::mutex>\{m\};} crée un temporaire, qui verrouille puis déverrouille au point-virgule. Donc \texttt{total += n} s'exécute sans protection, et il y a une course aux données. Il faut donner un nom au garde~: \texttt{std::lock\_guard<std::mutex> lock\{m\};}.

L'objet n'a pas de nom, donc sa durée de vie s'achève avec son expression complète, avant la ligne suivante. Le compilateur voit une expression légale, avec un constructeur et un destructeur qui ont des effets de bord, et n'émet aucun avertissement, même avec \texttt{-Wall -Wextra -Werror}. ThreadSanitizer signale aussitôt une \texttt{data race} sur \texttt{total}. Sans cela, le total peut être juste ou faux, selon l'entrelacement des threads. Avec le nom ajouté, TSan reste silencieux et le programme affiche 200000.

Avec des parenthèses, \texttt{std::lock\_guard<std::mutex>(m);} s'analyse comme la déclaration d'une variable nommée \texttt{m} de type \texttt{lock\_guard}, sans initialiseur. Les parenthèses supplémentaires autour d'un déclarateur sont autorisées. Comme \texttt{lock\_guard} n'a pas de constructeur par défaut, le code ne compile pas ici, par chance~: le test le confirme. Pour un type doté d'un constructeur par défaut, il compilerait et masquerait le \texttt{m} global.

\iqfollow{Est-ce que \texttt{std::scoped\_lock} rend cela plus sûr~?}{Pas à lui seul. \texttt{std::scoped\_lock lock;} sans aucun mutex compile, puisque le gabarit accepte un nombre quelconque de mutex, y compris aucun, et il ne verrouille rien. TSan signale toujours la course aux données sur cette version. Il faut nommer le garde, et nommer le mutex qu'il contient.}

\iqwhy{C'est un bogue d'un seul jeton que nul avertissement du compilateur ne détecte, et repérer un garde sans nom est un bon test pour vérifier qu'on a vraiment lu du code de verrouillage.}
\end{solution}

\begin{solution}{iq:iv3:reading-debugging-and-software-design:14}
\texttt{transfer(x, y)} verrouille \texttt{x} puis \texttt{y}, et \texttt{transfer(y, x)} verrouille \texttt{y} puis \texttt{x}. Si deux appels de ce type se chevauchent, chacun peut tenir un mutex et attendre l'autre indéfiniment. TSan le signale bien dans cette exécution, sous le nom de \texttt{lock-order-inversion}. Il faut utiliser \texttt{std::scoped\_lock lock(from.m, to.m);}.

Cette exécution ne peut pas se bloquer, car le premier thread est joint avant que le second ne démarre. C'est là tout l'enjeu de la question~: on ne peut pas repérer cela avec un test qui passe par hasard. TSan n'a pas besoin que l'interblocage se produise. Il construit un graphe indiquant quel mutex était tenu lorsqu'un autre était acquis, et signale un cycle. Ici, il affiche \texttt{Cycle in lock order graph}, sur une exécution qui a pourtant donné la bonne réponse.

\texttt{std::scoped\_lock} avec deux mutex utilise \texttt{std::lock}, dont la norme dit qu'elle n'aboutit pas à un interblocage. La technique habituelle consiste à verrouiller l'un puis à essayer les autres, en relâchant en cas d'échec. Une autre correction est un ordre global fixe, par exemple selon l'adresse ou l'identifiant du compte. Cela permet aussi de voir la règle dans le code, et n'impose aucune nouvelle tentative.

Le piège, c'est le transfert vers soi-même. Avec \texttt{from} et \texttt{to} désignant le même compte, \texttt{scoped\_lock} tente de verrouiller deux fois le même mutex, et un \texttt{std::mutex} verrouillé deux fois par un même thread a un comportement indéfini. Dans le test, le programme est resté bloqué jusqu'au délai d'expiration. Il faut donc ajouter une vérification en tête~: \texttt{if (\&from == \&to) return;}, ou la règle équivalente pour le domaine considéré.

\iqwrong{«~Toujours verrouiller dans l'ordre des arguments.~» C'est là le bogue. L'ordre doit être une propriété des mutex, et non de l'appel.}

\iqwhy{Deux verrous pris dans des ordres opposés sur des chemins de code différents forment l'interblocage classique. Une bonne réponse mentionne l'outil, la correction, et le cas limite que la correction introduit.}
\end{solution}

\begin{solution}{iq:iv3:reading-debugging-and-software-design:15}
Le bogue réside dans le \texttt{if}~: après le retour de \texttt{wait}, rien ne vérifie que la file contient encore un élément. Il faut utiliser la forme à prédicat, \texttt{cv.wait(lock, [] \{ return !q.empty(); \});}. \texttt{notify\_one} seul ne corrigerait pas le problème.

Voici comment cela échoue. Les deux consommateurs attendent. \texttt{push(7)} appelle \texttt{notify\_all}, donc les deux se réveillent et se disputent le mutex. Le premier prend 7. Le second obtient alors le mutex, sort de \texttt{wait}, et va directement à \texttt{q.front()} sur une deque vide. Cet appel a un comportement indéfini. Sans vérifications activées, rien n'arrête la lecture, ni le \texttt{pop\_front} qui suit, et cela peut passer inaperçu pendant longtemps.

Pour rendre l'échec visible, il faut compiler avec \texttt{-D\_GLIBCXX\_ASSERTIONS}, ce qui active les contrôles de préconditions de la bibliothèque. L'exécution du test s'interrompt avec SIGABRT sur \texttt{front}. Avec le prédicat, la même compilation affiche les deux valeurs, et TSan reste silencieux.

\texttt{notify\_one} réveille un seul thread, ce qui paraît suffisant. Cela laisse pourtant deux façons d'échouer. D'abord, \texttt{wait} peut revenir sans aucune notification~: la formulation de la norme inclut «~ou de manière parasite~». Ensuite, entre la notification et le moment où le thread réveillé obtient le mutex, un troisième thread qui appelle \texttt{pop} peut prendre l'élément, car il trouve la file non vide et n'attend pas.

La règle tient en peu de mots~: une variable de condition ne garantit jamais que la condition est vraie, elle indique seulement qu'il vaut la peine de la vérifier à nouveau. Il faut donc toujours attendre dans une boucle, ou utiliser la forme à prédicat, qui n'est autre que cette boucle.

\iqwrong{«~Le \texttt{if} convient, car \texttt{push} ajoute toujours un élément avant de notifier.~» C'est exact, mais l'élément peut avoir disparu au moment où ce thread détient le verrou.}

\iqwhy{Cela apparaît dans presque toutes les files que les candidats écrivent au tableau. Les recruteurs cherchent la boucle, puis demandent pourquoi.}
\end{solution}

\begin{solution}{iq:iv3:reading-debugging-and-software-design:16}
Il affiche \texttt{copies 1, moves 1}. Le paramètre est \texttt{const Tracked}, et \texttt{std::move} appliqué à un objet const donne une rvalue const, qui se lie au constructeur de copie. Si l'on retire le \texttt{const}, il affiche \texttt{copies 0, moves 2}.

\texttt{std::move} ne déplace rien. C'est un transtypage vers une référence rvalue, et il conserve le \texttt{const}. La résolution de surcharge a alors deux candidats~: \texttt{Tracked(Tracked\&\&)} exige un argument non-const, donc ne convient pas, alors que \texttt{Tracked(const Tracked\&)} convient. Le compilateur copie donc discrètement, sans avertissement. Le seul déplacement est le déplacement extérieur, de \texttt{t} vers le paramètre. La seule copie va du paramètre vers \texttt{item\_}.

Le \texttt{const} sur un paramètre passé par valeur semble inoffensif, voire vertueux, puisqu'il ne modifie pas la signature de la fonction. Pour un paramètre dont on destine le contenu à être déplacé, il est précisément une erreur. Il en va de même pour une variable locale \texttt{const} que l'on déplacera plus tard avec \texttt{std::move}.

Dans \texttt{std::vector}, la réallocation doit offrir la garantie forte contre les exceptions~: si quelque chose lève une exception pendant le déplacement des éléments, l'ancien tampon doit rester intact. Un déplacement qui lève une exception aurait détruit les anciens éléments. Le vecteur utilise donc \texttt{std::move\_if\_noexcept}~: il déplace seulement si le constructeur de déplacement est \texttt{noexcept}, ou si le type ne peut pas être copié. Le test le montre~: deux appels à \texttt{emplace\_back} affichent \texttt{copies 0, moves 1} avec \texttt{noexcept}, et \texttt{copies 1, moves 0} sans cette spécification.

\iqwrong{«~\texttt{std::move} force un déplacement.~» Elle ne fait que rendre un déplacement possible. Compter avec un type suivi, comme ici, permet de le vérifier.}

\iqwhy{Les copies non remarquées de gros messages dans un chemin critique (hot path) viennent exactement de ces deux lignes. Elles ne produisent ni erreur de compilation, ni échec de test, seulement un programme plus lent.}
\end{solution}

\begin{solution}{iq:iv3:reading-debugging-and-software-design:17}
Ce n'est pas correct~: le second \texttt{publish} appelle une lambda qui contient un \texttt{this} pendant, et il écrit dans une mémoire déjà libérée. ASan signale \texttt{heap-use-after-free} à \texttt{total\_ += v}. La correction donne à l'abonnement un propriétaire, de sorte que le destructeur le désabonne.

Le flux stocke le \texttt{std::function}, qui stocke le \texttt{this} brut. Le flux ignore tout de la durée de vie de la stratégie. Après \texttt{s.reset()}, la stratégie n'existe plus, mais sa lambda est toujours dans \texttt{subs\_}. Compilé normalement, l'écriture tombe sur une mémoire libérée et rien ne se voit~: on obtient \texttt{1}, \texttt{done}. En production, la même écriture tombe sur ce qui a été alloué ensuite à cet endroit, et le symptôme apparaît ailleurs.

La correction présentée dans le \cref{lst:iv3:rd:callbackfixed} fait renvoyer à \texttt{subscribe} un descripteur, stocké dans la stratégie, et le destructeur le rend. Une liste conserve la validité des itérateurs pendant que d'autres abonnés arrivent et partent. Le programme passe sans erreur sous ASan.

Il existe un second bogue à prévenir. Le constructeur de copie par défaut copierait \texttt{slot\_}, mais laisserait la lambda pointer vers l'objet d'origine, et les deux destructeurs effaceraient le même emplacement. La classe doit donc interdire la copie, comme ici. Le test montre que la copie ne compile alors pas.

Une autre conception courante utilise \texttt{std::shared\_ptr} pour la stratégie et capture un \texttt{std::weak\_ptr} dans la lambda. Le flux le verrouille à chaque appel, et un pointeur faible expiré est ignoré. C'est plus simple avec plusieurs threads, mais cela coûte une opération atomique par appel.

\omcode{code/interviews-3/09-reading-debugging-and-software-design/cpp/callback_fixed.cpp}{8}{45}{L'abonnement a un propriétaire.\protect\hfil\protect\break}\label{lst:iv3:rd:callbackfixed}

\iqfollow{Que se passe-t-il si un callback se désabonne lui-même pendant \texttt{publish}~?}{Cette version efface alors l'élément sur lequel se trouve la boucle range-for, ce qui est un comportement indéfini. La réponse usuelle consiste à le marquer comme mort et à le retirer après la boucle, ou à parcourir une copie.}

\iqwhy{Les callbacks sont le moyen par lequel les flux atteignent les stratégies. Ordonner correctement la destruction est un travail quotidien dans un système de trading.}
\end{solution}

\begin{solution}{iq:iv3:reading-debugging-and-software-design:18}
Il affiche \texttt{False 1 True}~: l'ordre est dans l'ensemble, mais \texttt{in} ne parvient pas à le trouver. Il faut calculer le hash sur \texttt{oid} seul, la partie qui ne change jamais, et faire en sorte que \texttt{\_\_eq\_\_} lui soit cohérent. Si l'on supprime \texttt{\_\_hash\_\_}, \texttt{book.add(o)} lève une \texttt{TypeError}.

Un ensemble range chaque élément dans un emplacement choisi d'après son hash au moment de l'insertion. Modifier \texttt{qty} change le hash que calcule désormais \texttt{\_\_hash\_\_}, si bien que la recherche regarde dans un autre emplacement. Elle ne trouve rien et ne compare jamais. L'élément est pourtant toujours là, comme le montrent les deux dernières valeurs~: \texttt{len} vaut 1, et un parcours par identité le trouve. On dispose alors d'un ordre qu'on ne peut ni trouver ni retirer.

Le modèle de données de Python met directement en garde contre cela~: une classe qui définit des objets mutables et implémente \texttt{\_\_eq\_\_} ne doit pas implémenter \texttt{\_\_hash\_\_}, car la valeur de hash d'une clé doit être immuable. Le langage n'en impose qu'une moitié. Une classe qui définit \texttt{\_\_eq\_\_} sans \texttt{\_\_hash\_\_} a son hash fixé à \texttt{None}, et ses instances ne sont donc pas hachables. Le test montre cette erreur.

Avec \texttt{hash(self.oid)} et l'égalité sur \texttt{oid}, la ligne affiche \texttt{True 1 True}. Mais cela a une conséquence~: deux ordres de même identifiant sont désormais égaux quelle que soit leur quantité. C'est ce que l'on veut pour l'identité dans un livre d'ordres. Ce n'est pas ce que l'on veut si l'on compare deux instantanés. Une conception plus propre sépare les deux~: un dictionnaire qui associe l'identifiant à l'ordre garde l'identité dans la clé et laisse l'ordre libre de changer.

\iqwrong{«~Les ensembles se comparent avec \texttt{==}, donc il le trouvera.~» Ils ne comparent qu'à l'intérieur d'un emplacement choisi par le hash, si bien que la comparaison n'a jamais lieu.}

\iqwhy{Hacher un champ mutable est une erreur facile au symptôme pénible~: des éléments bien présents et pourtant introuvables.}
\end{solution}

\begin{solution}{iq:iv3:reading-debugging-and-software-design:19}
Il affiche \texttt{150.0} et \texttt{False}~: après \texttt{del p} et une collecte complète du ramasse-miettes, le pricer est toujours vivant. Le cache le maintient en vie, avec sa courbe d'un million d'éléments. Il faut donner à chaque instance son propre cache.

\texttt{@lru\_cache} sur une méthode enveloppe la fonction une seule fois, au niveau de la classe. Sa clé contient tous les arguments, et le premier d'entre eux est \texttt{self}. Le cache de niveau classe détient donc une référence forte vers chaque instance qui l'a jamais appelé. La FAQ de Python dit la même chose de \texttt{lru\_cache} sur les méthodes~: «~il crée une référence vers l'instance~». Avec \texttt{maxsize=None}, il n'évince jamais rien, si bien que chaque nouveau \texttt{Pricer} et sa courbe restent en mémoire jusqu'à la fin du processus. C'est là la croissance mémoire. Les linteurs signalent aussi ce schéma~: la règle B019 de ruff signale précisément cette ligne.

Pour conserver la mise en cache, il faut créer le cache dans le constructeur~:

\texttt{self.price = lru\_cache(maxsize=None)(self.\_price)}

Le cache appartient alors à l'instance. Il se forme désormais un cycle de références (instance vers cache, vers méthode liée, vers instance), si bien que l'objet est libéré par le ramasse-miettes de cycles et non immédiatement. Le test affiche \texttt{True} après \texttt{gc.collect()}. Borner \texttt{maxsize} est aussi judicieux~: avec la forme au niveau de la classe, cela limite les dégâts sans y remédier, car quelques gros objets peuvent rester épinglés.

\iqfollow{Cette classe cache-t-elle encore un autre bogue~?}{Oui. La clé du cache n'inclut pas le contenu de \texttt{self.curve}. Si quelqu'un met à jour la courbe sur place, le cache renvoie des prix périmés. Un cache doit être indexé sur tout ce dont dépend le résultat, ou bien l'entrée doit être immuable.}

\iqwhy{La mise en cache dans un pricer est naturelle, et c'est la fuite mémoire réelle la plus courante dans la recherche et les services de pricing en Python.}
\end{solution}

\begin{solution}{iq:iv3:reading-debugging-and-software-design:20}
Il affiche \texttt{1 day, 0:00:00}, puis \texttt{23:00:00}. Le second est exact~: le 8 mars 2026, New York a avancé ses horloges, si bien que ce jour-là a compté 23 heures.

La documentation de Python énonce la règle~: si les deux datetime sont conscients du fuseau et ont le même attribut \texttt{tzinfo}, les attributs \texttt{tzinfo} sont ignorés et aucun ajustement de fuseau n'est effectué. Ici, les deux utilisent le même objet \texttt{ny}, donc la soustraction opère sur les champs de l'heure murale, de minuit à minuit, et donne un jour. Converties en UTC, les deux dates portent des décalages différents~: le test montre $-5$ heures juste avant le changement et $-4$ après. La soustraction donne alors le temps écoulé. Le même piège joue dans l'autre sens le 1 novembre 2026, jour de 25 heures.

Tout ce qui divise par un temps écoulé déraille silencieusement avec la première forme. Un volume par heure, ou une fourchette moyenne pondérée par le temps pour cette journée, est faussé d'environ 4~\%, puisque 24 heures sont retenues pour une journée de 23 heures. Un délai d'attente ou un calendrier de séance calculé en ajoutant des heures à une heure locale peut être décalé d'une heure. Une annualisation ou un temps jusqu'à l'échéance calculés sur des horodatages locaux peuvent subir la même erreur.

La règle dans les sociétés de trading~: stocker et calculer en UTC, ou en nanosecondes entières depuis l'époque, et ne convertir en heure locale que pour l'affichage, ou pour appliquer les règles de séance des bourses, qui sont elles-mêmes définies en heure locale.

\iqwrong{«~\texttt{zoneinfo} gère l'heure d'été.~» Il donne bien les bons décalages, mais la soustraction les ignore lorsque les deux membres ont le même \texttt{tzinfo}, comme ici.}

\iqwhy{Deux jours par an, cela représente une heure de données fausses, en avance ou en retard à l'ouverture. Le problème apparaît d'abord en production, car peu de tests utilisent ces dates.}
\end{solution}

\begin{solution}{iq:iv3:reading-debugging-and-software-design:21}
Il affiche \texttt{2 2} et rien de plus. On ne voit jamais \texttt{drop a} ni \texttt{drop b}, car les deux nœuds se possèdent mutuellement et aucun des deux compteurs n'atteint zéro. Selon les règles de Rust, ce n'est pas un bogue de sûreté. Il faut remplacer \texttt{Rc} par \texttt{Weak} dans le champ \texttt{next}.

Chaque nœud démarre avec une référence forte issue de sa variable locale, et l'autre nœud en ajoute une, d'où \texttt{2 2}. À la fin de \texttt{main}, les deux variables locales disparaissent et chaque compteur tombe à 1. Rien ne peut plus les faire tomber à zéro, puisque chacun est maintenu en vie par l'autre. La mémoire n'est donc jamais libérée et les destructeurs ne s'exécutent jamais.

La bibliothèque standard le dit clairement~: un cycle de pointeurs \texttt{Rc} ne sera jamais désalloué, et \texttt{Weak} existe précisément pour rompre les cycles. La Rust Reference range les fuites de mémoire, ainsi que les interblocages, parmi les comportements non considérés comme dangereux. Il ne faut donc jamais compter sur \texttt{Drop} pour libérer une limite ou clôturer une position.

Avec \texttt{RefCell<Weak<Node>{}>} et \texttt{Rc::downgrade}, la version corrigée affiche \texttt{1 1}, puis \texttt{a} (la récupération depuis \texttt{b} réussit toujours), puis \texttt{drop b}, \texttt{drop a}. Un \texttt{Weak} ne maintient pas la valeur en vie~: \texttt{upgrade} renvoie donc une \texttt{Option}, et il faut traiter le cas \texttt{None}.

\iqwrong{«~Rust empêche les fuites de mémoire.~» Il empêche l'utilisation après libération et les courses aux données dans le code sûr. Une fuite reste sans danger pour la mémoire.}

\iqfollow{Quel pointeur doit être faible dans un arbre~?}{Celui qui va de l'enfant vers le parent. Un parent possède ses enfants, et un enfant n'a besoin d'atteindre son parent que tant que celui-ci existe.}

\iqwhy{Un état en forme de graphe, comme des ordres qui pointent vers leurs niveaux de prix et des niveaux qui pointent vers leurs ordres, incite à former des cycles en Rust comme dans tout langage à comptage de références.}
\end{solution}

\begin{solution}{iq:iv3:reading-debugging-and-software-design:22}
Avec rustc 1.97, le programme se bloque indéfiniment. Le garde issu du premier \texttt{lock()} vit jusqu'à la fin du \texttt{match}, et la branche tente de reverrouiller le même mutex. L'édition 2024 ne change rien à cela. Il faut d'abord ajouter un \texttt{let}~: \texttt{let item = queue.lock().unwrap().pop();}, puis
\texttt{match item \{}. 
\texttt{queue.lock().unwrap()} crée un \texttt{MutexGuard} temporaire, et \texttt{pop()} n'emprunte à ce garde que brièvement. On pourrait penser que le garde disparaît dès que \texttt{pop} rend la main. Mais la Rust Reference précise que l'expression évaluée par un \texttt{match} n'est pas une portée temporaire, si bien que les temporaires de cette expression peuvent être détruits après l'expression \texttt{match}. Le garde reste vivant à travers les branches. Le second \texttt{lock()} sur le même thread attend alors un verrou que ce thread détient déjà. La documentation de \texttt{Mutex::lock} indique que cet appel peut paniquer. Ici il est resté bloqué, et le test constate le dépassement de délai~: il ne faut donc pas compter sur la panique.

L'édition 2024 modifie deux portées, et aucune des deux n'est celle-ci. Les temporaires de l'expression d'un \texttt{if let} sont désormais détruits avant la branche \texttt{else}, mais pas avant le corps, et les temporaires de l'expression finale sont détruits plus tôt. Le test construit ce \texttt{match} ainsi que la forme \texttt{if let} sous l'édition 2024~: les deux se bloquent encore, puisque le second verrou se trouve dans le corps.

Avec le \texttt{let}, l'instruction s'achève à son point-virgule, le garde est détruit, et le programme affiche \texttt{requeued 3}, puis \texttt{done}.

\iqwrong{«~L'analyseur d'emprunts empêche les interblocages.~» Il empêche les courses aux données. La Rust Reference range les interblocages parmi les comportements non considérés comme dangereux.}

\iqwhy{Rien dans le texte ne montre le garde, encore moins la durée de sa vie. Lire pour repérer cette portée est une compétence, et c'est celle qui permet de trouver cet interblocage avant la production.}
\end{solution}

\begin{solution}{iq:iv3:reading-debugging-and-software-design:23}
Il affiche \texttt{true}, puis \texttt{null}. \texttt{Key} redéfinit \texttt{equals} mais pas \texttt{hashCode}, donc deux clés égales ont des codes de hachage différents, et la table ne retrouve pas l'entrée. Il faut ajouter \texttt{hashCode}, qui renvoie \texttt{symbol.hashCode()}.

La documentation de \texttt{Object} énonce le contrat~: si deux objets sont égaux selon \texttt{equals}, appeler \texttt{hashCode} sur chacun doit produire le même entier, si bien qu'il est en général nécessaire de redéfinir \texttt{hashCode} dès que \texttt{equals} est redéfinie. Le \texttt{hashCode} par défaut ne repose pas sur les champs, si bien que des clés égales ont presque toujours des hachages différents. \texttt{HashMap.get} compare d'abord le hachage stocké, et n'appelle \texttt{equals} que s'ils concordent, si bien que même une clé située dans le même compartiment n'est pas trouvée. La recherche renvoie donc \texttt{null} sans jamais appeler \texttt{equals} de la classe. Avec la redéfinition, le programme affiche \texttt{true} et \texttt{1000}.

Le compilateur ne dit rien, car les deux méthodes sont facultatives. Un \texttt{null} renvoyé par \texttt{get} ressemble aussi à «~aucune limite définie~». C'est là le point dangereux. Dans un contrôle de limite, déballer ce \texttt{null} en \texttt{long} produit au mieux une \texttt{NullPointerException}. Au pire, il se lit comme «~pas de limite, donc autoriser~».

Autre façon de casser le mécanisme~: garder les deux méthodes, mais y utiliser un champ mutable, puis modifier ce champ après l'insertion. Le hachage diffère alors de celui stocké, et la même recherche renvoie \texttt{null}. C'est pourquoi les classes de clés doivent être immuables. Ici, \texttt{symbol} est \texttt{final}, et la classe est elle aussi \texttt{final}, si bien qu'une sous-classe ne peut pas modifier ce que signifie l'égalité. Le test exécute aussi cette variante mutable.

\iqwrong{«~Il affiche 1000, car les clés sont égales.~» \texttt{HashMap} n'atteint jamais le contrôle d'égalité.}

\iqwhy{C'est la première question d'un entretien Java dans une banque, car chaque table de positions et de limites repose dessus.}
\end{solution}

\begin{solution}{iq:iv3:reading-debugging-and-software-design:24}
Il affiche un nombre très inférieur à 400000, environ 130000 dans deux exécutions ici avec quatre threads, et ce nombre change d'une exécution à l'autre. La table est sûre thread par thread, mais \texttt{get} suivi de \texttt{put} constitue deux appels, si bien que des mises à jour se perdent entre les deux. Remplacer les deux lignes par \texttt{position.merge(symbol, qty, Long::sum);}.

Chaque opération de \texttt{ConcurrentHashMap} est sûre, et la structure interne de la table n'est jamais corrompue. Mais le thread A peut lire 5, le thread B peut lire 5, les deux calculent 6 et les deux écrivent 6. Un incrément est perdu. C'est une course de type «~vérifier puis agir~», et aucun choix de collection ne vous en protège tant que la logique reste répartie sur deux appels.

\texttt{merge} effectue la lecture, le calcul et l'écriture en une seule étape sur cette clé. La documentation indique que l'invocation entière de la méthode est réalisée de façon atomique, et que les mises à jour par d'autres threads peuvent être bloquées entre-temps. Garder donc la fonction courte, et ne pas toucher à la table depuis l'intérieur. Avec cette ligne, le test affiche exactement 400000.

\texttt{compute} fait la même chose avec une fonction plus générale, et \texttt{putIfAbsent} couvre le cas de la première fois. Ici, \texttt{Long::sum} suffit.

Pour un compteur sollicité par de nombreux threads et lu rarement, utiliser \texttt{LongAdder}. Il conserve une ou plusieurs variables dont le nombre augmente en cas de contention, et il les additionne à la lecture. Sa Javadoc indique qu'il est en général préférable à \texttt{AtomicLong} pour de tels cumuls, au prix d'un peu plus de mémoire. \texttt{map.computeIfAbsent(symbol, s -> new LongAdder()).add(qty)} en est la forme habituelle.

\iqwrong{«~Passez la table en \texttt{Hashtable} ou enveloppez-la dans \texttt{synchronizedMap}.~» Ces structures verrouillent chaque appel, si bien que la course entre les deux appels subsiste.}

\iqwhy{La tenue des positions dans un service concurrent est exactement ce code. Un nombre qui dérive vers le bas compte parmi les bogues les plus graves, car rien ne lève d'exception.}
\end{solution}

\begin{solution}{iq:iv3:reading-debugging-and-software-design:25}
Il affiche \texttt{True}~: la stratégie survit à la collecte. L'événement de \texttt{feed} contient un délégué qui contient la stratégie. Des stratégies de courte durée s'accumuleraient donc, chacune recevant encore les cotations. Désabonnez-vous avec \texttt{-=}, normalement dans \texttt{Dispose}.

\texttt{feed.Tick += OnTick} crée un délégué dont la cible est l'instance de \texttt{Strategy}, et l'ajoute à la liste du flux. Le ramasse-miettes suit les références, si bien qu'un \texttt{feed} encore vivant maintient le délégué en vie, ce qui maintient la stratégie en vie. Le guide de programmation de Microsoft le dit presque dans les mêmes termes~: tant que l'objet éditeur conserve cette référence, le ramasse-miettes ne supprimera pas votre objet abonné, et vous devez donc vous désabonner avant de libérer l'abonné.

Le listing est construit avec soin pour le montrer honnêtement. \texttt{NoInlining} sur \texttt{Make} et la compilation en Release garantissent qu'aucune variable locale ne maintient la stratégie en vie par accident. \texttt{GC.KeepAlive(feed)} maintient le flux en vie jusqu'au contrôle. Il ne reste donc que le chemin par l'événement.

Des abonnés de courte durée à un éditeur de longue durée constituent la mauvaise configuration~: par exemple, la vue d'un instrument qui s'abonne à un flux global pendant toute la durée d'un écran. La mémoire augmente, et le processeur aussi, car chaque stratégie morte exécute encore \texttt{OnTick} à chaque cotation, en travaillant pour rien. Ce second symptôme est souvent celui que l'on remarque en premier.

Avec \texttt{IDisposable} et \texttt{feed.Tick -= OnTick} dans \texttt{Dispose}, ainsi que \texttt{using var s} dans \texttt{Make}, le programme affiche \texttt{False}.

\iqwrong{«~C\# dispose d'un ramasse-miettes, donc il ne peut pas y avoir de fuite.~» Le ramasse-miettes libère ce qui est inaccessible. Ici l'objet reste accessible, et c'est une fuite à tous égards.}

\iqfollow{La stratégie qui détient une référence vers le flux pose-t-elle le même problème~?}{Non. Cette référence va dans l'autre sens~: une stratégie par ailleurs inaccessible peut être collectée avec sa référence. Seule la relation de l'éditeur vers l'abonné maintient un abonné en vie.}

\iqwhy{Les applications de bureau et de valorisation de longue durée dans les banques rencontrent souvent ce problème.}
\end{solution}

\begin{solution}{iq:iv3:reading-debugging-and-software-design:26}
Seule la première version peut échouer sur un mauvais nombre, et elle ne peut pas détecter le bogue. \texttt{vwap} ignore les quantités et prend une moyenne simple. Tester avec des quantités inégales, et utiliser \texttt{pytest.raises} ainsi qu'un \texttt{assert} nu.

Reprenons-les. \texttt{test\_two\_trades} utilise des quantités égales, 50 et 50, et pour celles-ci la moyenne simple et la moyenne pondérée coïncident~: 101 dans les deux cas. \texttt{test\_empty} ne contient aucun \texttt{assert} et avale l'exception qu'il espère, si bien qu'il réussit si l'appel lève cette exception et réussit aussi si l'appel renvoie une valeur. \texttt{test\_one\_trade} affirme un tuple à deux éléments. Un tuple non vide est toujours vrai, si bien qu'il ne peut jamais échouer, même s'il compare à une valeur fausse, 99. Python lui-même avertit à la compilation~: \texttt{assertion is always true}. Le test capture cet avertissement. pytest comptabilise néanmoins \texttt{3 passed}.

La bonne fonction est la somme de $p \times q$ divisée par la somme de $q$. Réécrire~:

\texttt{assert vwap([(100.0, 10), (102.0, 30)]) == 101.5}

Avec la fonction actuelle, cela donne 101, et le test échoue. Calculez votre valeur attendue à la main, et ne la recopiez pas depuis une exécution.

\texttt{with pytest.raises(ValueError):} suivi de \texttt{vwap([])}, qui échoue si rien n'est levé. Puis \texttt{assert vwap([(100.0, 7)]) == 100.0}.

Ajouter des propriétés, qui détectent des familles entières de bogues~: le résultat se situe entre le prix le plus bas et le prix le plus haut, et multiplier toutes les quantités par 3 ne change rien. Le test les exécute toutes deux sur 200 listes aléatoires.

Pour vérifier les tests eux-mêmes, casser volontairement le code et constater qu'un test échoue. C'est l'idée qui sous-tend le test de mutation\index{test de mutation}, et ici il aurait révélé les trois défauts à la fois.

\iqwrong{«~La couverture est de 100~\%, donc c'est testé.~» Chaque ligne s'exécute. Rien n'est vérifié.}

\iqwhy{Les entreprises lisent vos tests aussi attentivement que votre code. Des tests qui ne peuvent pas échouer sont pires que l'absence de tests, car ils donnent une fausse confiance.}
\end{solution}

\begin{solution}{iq:iv3:reading-debugging-and-software-design:27}
Donner à chaque quantité son propre type~: \texttt{struct Price \{ long ticks; \};} et \texttt{struct Qty \{ long lots; \};}. \texttt{send(qty, price)} ne compile alors pas. Dans le code machine, n'attendre aucune différence avec un \texttt{long} simple.

Le \Cref{lst:iv3:rd:strong} est le résultat. Il compile avec les mêmes options, affiche la bonne valeur, et l'inversion des arguments devient une erreur~: le test le confirme. Pour le coût, on compare deux fonctions~:

\texttt{long cost\_raw(long, long)} et \texttt{long cost\_typed(Price, Qty)}.

Avec g++ 11.4 en \texttt{-O2} sur x86-64 Linux, les deux compilent vers les mêmes instructions, \texttt{movq}, \texttt{imulq}, \texttt{ret} (plus un \texttt{endbr64} avec les options par défaut d'Ubuntu). Le test les compare ligne à ligne. La raison tient à la convention d'appel System V~: une petite structure contenant un seul entier passe dans un registre exactement comme l'entier. Un enveloppement n'a aucune existence à l'exécution, tout se passe dans le vérificateur de types.

N'ajouter que les opérations qui ont un sens. Un prix moins un prix est un écart de prix, et une quantité multipliée par un prix est un notionnel. Un prix plus un prix n'a pas de sens~: l'ometter, et le compilateur le refuse désormais aussi. Les conversions depuis \texttt{long} doivent être explicites, sinon la protection disparaît.

Ce listing comporte une faille. Ce sont des agrégats, si bien que \texttt{send(\{500\}, \{101\})} compile encore et inverse toujours les arguments. Si cela vous préoccupe, ajoutez un constructeur \texttt{explicit}.

\omcode{code/interviews-3/09-reading-debugging-and-software-design/cpp/strong_types.cpp}{3}{18}{Le même ordre avec des types.\protect\hfil\protect\break}\label{lst:iv3:rd:strong}

\iqwrong{«~Commentez les paramètres, ou utilisez des variables nommées.~» Le listing nomme déjà ses variables, et c'est précisément là qu'il échoue.}

\iqwhy{Un prix et une quantité inversés constituent un bogue de saisie commis par un programmeur. Les enveloppes de type comptent parmi les rares outils qui suppriment une classe de bogues sans coût à l'exécution.}
\end{solution}

\begin{solution}{iq:iv3:reading-debugging-and-software-design:28}
La bisection simple demande 10 exécutions, $2^{10} = 1024$, soit 60 minutes. Avec le test instable, il faut 24 exécutions par commit, et le pire cas dure 24 heures.

La commande \texttt{git bisect} est une recherche dichotomique~: chaque test divise l'intervalle par deux, si bien que 1024 commits demandent 10 étapes. Cela suppose un commit bon, un commit ultérieur mauvais, et un seul changement de bon à mauvais entre les deux.

Avec le test instable, un échec est concluant, puisqu'un commit sans le bogue n'échoue jamais. Un succès ne l'est pas. Un mauvais commit réussit $r$ exécutions d'affilée avec une probabilité de $(3/4)^r$. Dans ce cas, la recherche prend un mauvais commit pour un bon, et part dans la mauvaise direction. Sur le pire chemin, les 10 commits testés sont tous mauvais, et il faut qu'ils soient tous correctement jugés~:
\[
\left(1 - (3/4)^r\right)^{10} \ge 0.99.
\]
En prenant la racine dixième, $(3/4)^r \le 1 - 0.99^{1/10} \approx 0.001$, donc $r \ge \ln(0.001) / \ln(0.75) \approx 24$. En effet $r = 24$ donne 0.9900 et $r = 23$ donne 0.987. Chaque mauvais commit exige donc un taux de manquement d'environ $1/1000$.

Pour le temps, le pire cas compte 10 commits à 24 exécutions chacun~: 240 exécutions, soit 24 heures. Chaque commit bon coûte les 24 exécutions, car il doit toutes les réussir. Un mauvais commit s'arrête à son premier échec, après environ 4 exécutions en moyenne. Un chemin typique rencontre 5 de chaque, donc $5 \times 24 + 5 \times 4 = 140$ exécutions, soit 14 heures. Une simulation de 4000 recherches dans le test trouve le bon commit plus de 99~\% du temps.

En pratique, rendre d'abord le test moins instable, ou l'exécuter sur de nombreuses machines à la fois. Avec 24 exécutions en parallèle, chaque étape dure 6 minutes et la recherche entière une heure.

\iqwrong{«~Relancez-le trois ou quatre fois, puisqu'il échoue une fois sur quatre.~» Quatre exécutions manquent encore le bogue avec une probabilité de 32~\%, et sur dix étapes la recherche a plus de chances de se tromper que de réussir.}

\iqwhy{Trouver à quel moment une panne rare est apparue est une tâche courante dans une entreprise, et l'arithmétique décide s'il s'agit d'un travail d'un après-midi ou d'une semaine perdue.}
\end{solution}

\begin{solution}{iq:iv3:reading-debugging-and-software-design:29}
Je demanderais une modification. \texttt{ReplayBook} rompt ce que promet \texttt{Book}, car tout code qui reçoit un \texttt{Book} peut appeler \texttt{add}, et ici cela lève une exception. Scinder l'interface~: une vue en lecture seule avec \texttt{total}, et un \texttt{Book} modifiable construit par-dessus.

Le test de l'héritage consiste à se demander si la sous-classe peut remplacer la classe de base partout où celle-ci est utilisée. C'est le principe de substitution de Liskov\index{principe de substitution de Liskov}, et \texttt{top\_up} montre la défaillance~: avec un \texttt{Book} il renvoie un total, et avec un \texttt{ReplayBook} il lève une exception. Rien dans les types n'avertit l'auteur de \texttt{top\_up}, si bien que la défaillance apparaît à l'exécution, probablement sur le chemin que personne ne teste, et seulement pour les livres rejoués.

Le sens de la hiérarchie est lui aussi erroné. La lecture seule est l'interface la plus petite, elle doit donc être la base. \texttt{ReadableBook} possède \texttt{total}~; \texttt{Book} ajoute \texttt{add}. \texttt{ReplayBook} est alors un \texttt{ReadableBook}, et \texttt{top\_up} déclare qu'il a besoin d'un \texttt{Book}. Avec un vérificateur statique ou un compilateur C++, lui passer un livre rejoué est une erreur au point même de la faute. En Python, on exprimerait cela avec \texttt{typing.Protocol} et on laisserait un vérificateur de types le faire respecter.

Souvent, il faut aller plus loin et renoncer à l'héritage. La relecture a besoin d'un chargeur qui remplit un livre, puis remet une vue en lecture seule. La composition offre cela sans créer aucune nouvelle classe de livres, et la même implémentation de \texttt{Book} reste la seule à tester.

Il existe un cas pour accepter la sous-classe~: un script de courte durée, un seul appelant, sous une échéance. Le dire, et laisser une note. J'accepterais aussi si le contrat de la classe de base déclarait déjà l'opération facultative, et si les appelants étaient écrits pour s'y attendre. Mais alors c'est ce contrat qu'il faut remettre en question.

\iqwrong{«~C'est acceptable, car \texttt{NotImplementedError} indique exactement ce qui s'est mal passé.~» Elle le signale à l'exécution, sur le chemin de relecture. Un type le signalerait avant toute exécution.}

\iqwhy{Les recruteurs veulent entendre une raison et un compromis, non le nom d'un principe. Une réponse de conception qui dit quand on ne s'en embarrasserait pas est plus convaincante.}
\end{solution}

\begin{solution}{iq:iv3:reading-debugging-and-software-design:30}
\texttt{new Stats} laisse \texttt{count} et \texttt{sum} non initialisés, et les lire est un comportement indéfini. Les tests passent parce que la mémoire fraîche du tas vaut par hasard zéro. Lancer avec \texttt{MALLOC\_PERTURB\_=165}~: le programme affiche alors 1. Écrire \texttt{new Stats\{\}}.

\texttt{Stats} n'a pas de constructeur, donc \texttt{new Stats} effectue une initialisation par défaut, qui pour une struct simple ne fait rien. \texttt{new Stats\{\}} effectue une initialisation à valeur, qui remet tout à zéro. Les deux diffèrent de deux caractères, et aucun des deux ne déclenche d'avertissement avec \texttt{-Wall -Wextra -Werror}, car l'allocation et l'utilisation se trouvent dans des fonctions différentes.

Pourquoi le test passe-t-il~? Dans un processus de courte durée, l'allocateur fournit de la mémoire fraîche venant du système d'exploitation, et celle-ci est à zéro. Dans un processus de longue durée, l'allocateur réutilise des blocs libérés, si bien qu'on voit d'anciens contenus, ce qui produit les absurdités observées en production. Le test unitaire et la production n'ont donc vraiment pas les mêmes conditions initiales.

glibc offre un moyen de le révéler sans recompiler. La page de manuel de \texttt{mallopt} décrit \texttt{M\_PERTURB}~: les octets de la mémoire allouée sont initialisés au complément de l'octet de poids faible de la valeur. Définir la variable d'environnement \texttt{MALLOC\_PERTURB\_=165} remplit les allocations provenant de l'allocateur principal avec \texttt{0x5a}, et le programme affiche alors 1~: les deux champs partent de la même très grande valeur, et l'ajout de 600 et de 3 ne modifie guère le rapport. Avec \texttt{new Stats\{\}}, la même exécution affiche 200. Sous glibc 2.35, un bloc réutilisé depuis le cache par thread conserve son contenu~; pour tous les attraper, définir aussi \texttt{GLIBC\_TUNABLES=glibc.malloc.tcache\_count=0}. La vérification ne demande aucune recompilation, elle peut donc tourner dans chaque tâche de test.

D'autres outils existent. g++ 11.4 refuse \texttt{-fsanitize=memory}, comme le confirme le test. Le MemorySanitizer de Clang est l'outil conçu pour cela, et il exige une recompilation avec Clang. Memcheck, de Valgrind, signale lui aussi l'usage de valeurs non initialisées. Aucun des deux n'a été exécuté pour cette réponse.

\iqwrong{«~Il suffit de mettre les champs à zéro dans un constructeur, partout.~» Cela fonctionne, mais les initialiseurs de membres par défaut, \texttt{long count = 0;}, sont plus courts et ne risquent pas d'être oubliés dans un nouveau constructeur.}

\iqwhy{«~Ça marche sur ma machine et dans les tests~» est exactement ce que produit ce bogue, et un développeur capable de le faire reproduire en une seule commande est celui qu'on veut avoir dans l'équipe.}
\end{solution}

\begin{solution}{iq:iv3:reading-debugging-and-software-design:31}
La mémoire reste atteignable~: \texttt{done} est une variable globale qui pointe vers chaque ordre. LeakSanitizer ne signale que la mémoire qui n'a plus aucun pointeur vivant à la sortie, il ne dit donc rien ici. Les trois causes possibles sont une vraie fuite, une croissance non bornée des données vivantes, et la rétention ou la fragmentation de l'allocateur.

Le test montre la différence. Le listing, exécuté sous ASan avec 100000 ordres, ne signale rien. Il suffit de remplacer l'insertion dans la map par un simple \texttt{new Order} dont on abandonne le pointeur~: l'outil signale alors la fuite immédiatement. L'outil répond à une question précise, «~quelque chose est-il inaccessible~?~», et ici sa réponse est juste. La faute est dans la conception~: c'est ce qu'on appelle une fuite logique.

Pour les distinguer, il faut procéder dans cet ordre. D'abord, relever dans le temps la vue qu'a l'allocateur de sa propre mémoire, à côté de la taille résidente. La fonction \texttt{mallinfo2} de glibc donne l'espace alloué et l'espace libre. Pour les octets alloués, additionner \texttt{uordblks} et \texttt{hblkhd}, car les grands blocs sont mappés directement et comptés à part.

Si les octets alloués croissent avec la mémoire résidente, il y a une croissance réelle de l'usage~: une fuite ou une structure qui ne fait que grandir. Un profileur de tas donne ensuite les sites d'allocation, classés par octets encore vivants. Un site dont la courbe ne cesse de monter est le suspect, et l'on lit ce qui le retient. Si LeakSanitizer reste muet, c'est que la mémoire est atteignable~: il faut donc chercher les maps, caches, files et journaux sans bornes.

Si les octets alloués restent stables et que la mémoire résidente croît, deux cas se présentent. Soit l'espace libre à l'intérieur de l'allocateur croît, soit la mémoire n'est jamais passée par \texttt{malloc}~: les mappages directs, les piles de threads et les bibliothèques qui ont leur propre allocateur n'apparaissent pas du tout, il faut donc vérifier \texttt{/proc/PID/smaps}. Dans le premier cas, l'allocateur retient des pages libérées qu'il ne peut pas rendre, généralement à cause d'un mélange de tailles qui fragmente ses arènes, ou de nombreux threads ayant chacun leur arène. Il faut regarder la valeur de l'espace libre, puis tester avec un autre allocateur ou une arène bornée.

Pour ce code~: chaque entrée coûte 136~octets de charge utile (8 pour la clé et 128 pour l'ordre), plus le surcoût du nœud d'arbre. Deux millions d'exécutions par jour représentent au moins 272~Mo par jour, ce qui correspond à la croissance observée. Il faut écrire l'audit sur disque, ne garder en mémoire que les enregistrements de la session, et ajouter une métrique sur la taille de la map.

\iqwrong{«~ASan ne signale aucune fuite, donc il s'agit de fragmentation.~» Cela ignore la cause la plus vraisemblable. Une croissance atteignable est bien plus fréquente.}

\iqwhy{Les processus de trading tournent pendant des jours sans redémarrage, et un problème de croissance est une investigation courante. Les recruteurs veulent une méthode ordonnée, non une liste d'outils.}
\end{solution}

\begin{solution}{iq:iv3:reading-debugging-and-software-design:32}
Ce n'est pas correct. La première lecture de \texttt{instance} n'est pas synchronisée avec l'écriture faite dans un autre thread~: il s'agit donc d'une course aux données et d'un comportement indéfini. x86-64 n'y change rien. La version simple et correcte est une variable statique locale à la fonction. Pour conserver le double contrôle, il faut \texttt{std::atomic} avec acquire et release.

L'échec classique~: le thread A exécute \texttt{new Config}, qui comporte trois étapes~: allouer, construire, affecter. Rien n'empêche l'affectation de devenir visible pour le thread B avant la construction~; B voit donc un pointeur non nul et lit un \texttt{limit} qui n'est pas encore fixé. B ne prend jamais le verrou, si bien que le verrou ne peut rien y faire.

Le modèle mémoire matériel fort de x86-64 garde les écritures dans l'ordre, mais les règles du langage s'imposent au compilateur, qui a le droit de réordonner une écriture non atomique, ou de supposer qu'aucun autre thread ne regarde. Et le programme est indéfini quoi que fasse la machine. TSan signale \texttt{data race} sur ce listing dans le test.

La correction la plus simple, également tirée du test~:

\texttt{Config* get() \{ static Config instance; return \&instance; \}}

La norme le garantit~: si plusieurs threads atteignent la déclaration pendant que la variable s'initialise, les autres attendent la fin de l'initialisation. g++ procède avec une variable de garde vérifiée à chaque appel, et \texttt{\_\_cxa\_guard\_acquire} sur le chemin lent~; c'est donc déjà un double contrôle, écrit d'avance par le compilateur. TSan ne signale rien.

Pour garder ce schéma, il faut utiliser \texttt{std::atomic<Config*>}. Le chemin rapide fait une lecture acquire. Le chemin lent tient le mutex, relit avec l'ordre relaxed, construit, puis fait une écriture release. Ainsi, un thread qui voit le pointeur voit aussi l'objet, et TSan ne signale rien sur cette version non plus. \texttt{std::call\_once} en est la forme fournie par la bibliothèque.

En Java, le même schéma fonctionne si le champ est \texttt{volatile}~: la Java Language Specification dit qu'une écriture dans un champ volatile précède au sens de happens-before chaque lecture ultérieure de ce champ. Sans \texttt{volatile}, le schéma est cassé de la même façon. Cela repose sur la spécification, car aucun test ici ne peut forcer le réordonnancement sur une JVM.

\iqwrong{«~x86 a un modèle mémoire fort, donc ça marche en pratique.~» C'est le compilateur qui réordonne ici, et un test qui passe prouve peu de chose sur une course aux données.}

\iqwhy{Ce schéma figure encore dans du vieux code C++ et Java, et il distingue ceux qui se souviennent de la recette de ceux qui savent dire ce qui ne va pas.}
\end{solution}

\begin{solution}{iq:iv3:reading-debugging-and-software-design:33}
Avec les fichiers objets dans l'ordre main, rates, limits, le programme affiche 1000. Avec main, limits, rates, il affiche 0. C'est le problème de l'ordre d'initialisation statique\index{ordre d'initialisation statique}. La solution consiste à utiliser une statique locale à la fonction, et à marquer une variable globale \texttt{constinit} pour en faire une erreur de compilation.

Les deux variables globales ont besoin d'un appel ou d'une lecture au démarrage, elles ont donc une initialisation dynamique. Dans un même fichier, l'ordre est celui des définitions. D'un fichier à l'autre, la norme laisse l'ordre indéterminé, il dépend donc de l'implémentation~; ici, avec g++ et GNU ld, il a suivi l'ordre de la ligne de commande de l'éditeur de liens.

Avant que toute initialisation dynamique ne s'exécute, toutes les variables globales sont mises à zéro. Ainsi, quand \texttt{limits} s'exécute en premier, il lit \texttt{base\_rate} comme 0 et calcule $0 \times 10$. Rien n'est non initialisé, aucun sanitizer ne se plaint, et la limite de risque vaut simplement zéro. Plus tard, \texttt{base\_rate} vaut 100, comme il se doit, mais la mauvaise valeur reste dans \texttt{limit}. Cela peut empirer~: une modification dans un script de construction, loin des deux fichiers, peut inverser l'ordre.

La correction transforme la variable globale en fonction~:

\texttt{long\& base\_rate() \{ static long v = compute\_base(); return v; \}}

Le premier appelant déclenche l'initialisation, quel que soit le fichier où il se trouve, et le test affiche 1000 pour les deux ordres d'édition des liens.

\texttt{constinit} est l'autre outil, introduit en C++20. Il exige qu'une variable soit initialisée à la compilation, et la norme rend le programme mal formé s'il a une initialisation dynamique. \texttt{constinit long limit = base\_rate * 10;} ne compile pas, ce qui est voulu~: il transforme un pari silencieux sur l'ordre en erreur. Il ne corrige rien à lui seul. Il indique quelles variables réécrire.

La conception la plus robuste n'a pas de telles variables globales. Il faut lire la configuration une seule fois dans \texttt{main}, construire les objets dans l'ordre voulu, et les passer là où ils sont nécessaires. L'ordre devient alors une ligne de code qu'on peut lire.

\iqwrong{«~Placer \texttt{limits.o} après \texttt{rates.o} et laisser un commentaire.~» Cela fonctionne avec cette chaîne d'outils, mais la norme ne promet rien. Une bibliothèque statique ou un nouvel éditeur de liens peut tout changer.}

\iqfollow{La version à statique locale est-elle sûre avec les threads~?}{Oui, l'initialisation est sûre vis-à-vis des threads depuis C++11. Il reste un piège~: l'ordre de destruction à la sortie, si le destructeur d'une autre variable globale l'utilise après qu'elle a disparu. Ici c'est un \texttt{long}, donc cela ne se produit pas.}

\iqwhy{Les limites et la configuration conservées dans des variables globales existent bel et bien, et une limite nulle trouvée en production, c'est une très mauvaise journée.}
\end{solution}

\begin{solution}{iq:iv3:reading-debugging-and-software-design:34}
\texttt{by\_virtual} et \texttt{by\_variant} contiennent un appel ou un saut indirect. \texttt{by\_final} et \texttt{by\_template} sont insérés en ligne en deux instructions (plus un \texttt{endbr64} avec les options par défaut d'Ubuntu), et \texttt{by\_switch} ne fait que des comparaisons et des branchements, sans appel indirect. La surprise vient de \texttt{by\_variant}~: beaucoup supposent que \texttt{std::visit} évite un appel indirect.

\texttt{by\_virtual} charge la table virtuelle et aboutit à \texttt{jmp *}~: le compilateur ne peut pas connaître le type. \texttt{by\_final} prend un \texttt{Momentum}, qui est \texttt{final}, donc il n'y a aucune redéfinition à chercher, et l'appel est inséré en ligne. \texttt{by\_template} connaît le type à la compilation, et le test trouve un code identique à celui de \texttt{by\_final}.

Sur libstdc++ 11, \texttt{std::visit} construit une table de pointeurs de fonctions indexée par l'alternative, puis utilise \texttt{call *(\%rdx,\%rax,8)}. Le test lit cela dans l'assembleur. Dans cette version, c'est donc un appel virtuel déguisé, auquel s'ajoute une vérification de l'état sans valeur. \texttt{by\_switch} teste l'indice avec \texttt{std::get\_if}, et tout s'insère en ligne en deux comparaisons et deux branchements, avec un chemin froid (cold path) qui lève une exception. Une bibliothèque ultérieure peut compiler \texttt{std::visit} autrement~; il faut donc lire la sortie de son propre compilateur, comme le fait ce test, avant de croire l'une ou l'autre affirmation.

Pour la vitesse, l'instruction elle-même coûte rarement le plus cher. La perte la plus importante tient à ce que le compilateur ne peut pas insérer en ligne à travers un appel indirect, ni optimiser au-delà de celui-ci. Il faut la mesurer sur place.

D'autres éléments décident généralement de la conception~:

Qui peut ajouter un type~? Si des stratégies arrivent d'autres équipes et de plug-ins, l'ensemble est ouvert, et les fonctions virtuelles sont le bon outil. Ajouter un type ne touche alors aucun code existant. Si l'ensemble est fermé et connu, un variant permet au compilateur de vérifier que chaque opération traite chaque type.

Où vit-il~? Un variant a la taille de sa plus grande alternative, stocké sur place sans allocation. Cela convient à un vecteur de messages. Les objets virtuels demandent surtout le tas, et le parcours de pointeurs coûte plus cher que l'appel.

Les templates donnent le meilleur code lorsque le type est fixé à la compilation. Mais ils se répandent dans les en-têtes, ralentissent la compilation et ne peuvent pas être choisis par un fichier de configuration à l'exécution.

\iqwrong{«~Les variants sont toujours plus rapides que les appels virtuels.~» Comme le montre ce compilateur, ce n'est même pas un autre genre d'appel.}

\iqwhy{Les entretiens sur la latence faible posent cette question pour voir si l'on se tourne vers l'assembleur, et si l'on sait argumenter la conception au-delà de l'ordre de la nanoseconde.}
\end{solution}

\begin{solution}{iq:iv3:reading-debugging-and-software-design:35}
Le programme affiche 40. Il devrait afficher 70, car la retransmission est la même exécution. Il faut rendre \texttt{on\_fill} idempotent en retenant le plus haut numéro de séquence appliqué par session, et en ignorant tout ce qui est égal ou inférieur à lui. Un écart dans la séquence doit arrêter le keeper.

La faute vient de ce que \texttt{on\_fill} traite chaque message comme un événement nouveau. Sur toute connexion qui renvoie des messages après une reconnexion, la livraison est au moins une fois, et non exactement une fois~; le récepteur doit donc rendre les répétitions sans effet. Le \Cref{lst:iv3:rd:replayfixed} en donne la forme.

\enlargethispage{-3\baselineskip}%
\omcode{code/interviews-3/09-reading-debugging-and-software-design/python/rd_replay_fixed.py}{1}{12}{Un keeper qui ignore les répétitions et s'arrête sur les écarts de séquence.\protect\hfil\protect\break}\label{lst:iv3:rd:replayfixed}

Trois décisions comptent ici. D'abord, la séquence est propre à chaque session~: l'état est donc une table, et le test vérifie que deux sessions ne se perturbent pas. Ensuite, un écart lève une erreur. Appliquer silencieusement l'exécution 4 après l'exécution 2 masque une exécution manquante, et la position devient alors fausse sans que personne ne le remarque. Un système de trading doit cesser de coter et demander la plage manquante. Enfin, le contrôle et la mise à jour vont ensemble~: l'état ne change que lorsqu'une exécution est appliquée.

Un high-water mark suppose que les exécutions arrivent dans l'ordre sur chaque session, et il ne stocke qu'un seul nombre. Si l'ordre n'est pas garanti, ou si les numéros de séquence repartent de zéro, il faut plutôt indexer sur un identifiant d'exécution unique et retenir les identifiants déjà vus. Cet ensemble grossit, il faut donc le tronquer, par exemple à la fin de la session. Quelle que soit la solution retenue, le high-water mark doit survivre à un redémarrage avec la position. Il faut les enregistrer ensemble, en une seule écriture atomique, ou reconstruire les deux à partir du journal. Si l'on sauvegarde la position sans le high-water mark, la première retransmission après un redémarrage est de nouveau comptée deux fois.

Pour les tests, prendre une suite aléatoire d'exécutions et la livrer avec des retransmissions aléatoires des derniers messages, entre les nouveaux. La position finale doit être égale à la somme simple des exécutions, quel que soit le schéma de retransmission. Le test est répété 200 fois. Ajouter un test pour un écart, et un test qui redémarre au milieu de la séquence.

\iqwrong{«~Réconcilier avec la bourse en fin de journée.~» C'est nécessaire, mais une position fausse pendant des heures peut dépasser une limite entre-temps, dans un sens comme dans l'autre.}

\iqwhy{Les reconnexions arrivent chaque jour, et une exécution comptée deux fois fait trader sur une position fausse. C'est l'une des fautes réelles les plus fréquentes dans un keeper de position.}
\end{solution}

\begin{solution}{iq:iv3:reading-debugging-and-software-design:36}
Il envoie un seul ordre. Il devrait en envoyer cinq, puisque les ticks sont espacés de 2 secondes en temps de marché. La stratégie lit l'horloge de la machine. Il faut injecter l'horloge et la passerelle, pour que la stratégie ne découvre jamais dans quel monde elle se trouve.

\texttt{time.time()} est une entrée cachée. Pendant un rejeu, les cinq ticks arrivent en quelques microsecondes de temps mural, de sorte que la limitation à 1 seconde bloque les quatre autres. En production, le même code envoie les cinq. Un backtest de ce genre serait faux d'une manière difficile à voir, car il tourne sans erreur et donne des chiffres. Il dépendrait aussi de la vitesse de la machine.

Le \Cref{lst:iv3:rd:strategyfixed} reprend la même logique, à ceci près qu'il interroge \texttt{self.clock.now()} et qu'il envoie par la passerelle qu'on lui donne. Dans un backtest, la boucle de rejeu règle l'horloge simulée sur l'horodatage de chaque événement avant de le livrer. En production, l'horloge enveloppe l'horloge système et la passerelle envoie de vrais ordres. Le test rejoue les mêmes ticks et obtient les cinq ordres, deux fois de suite, à l'identique.

\omcode{code/interviews-3/09-reading-debugging-and-software-design/python/rd_strategy_fixed.py}{1}{22}{L'horloge et la passerelle lui sont fournies.\protect\hfil\protect\break}\label{lst:iv3:rd:strategyfixed}

On note aussi la valeur de départ de la version corrigée~: \texttt{self.last} vaut moins l'infini. Dans l'original, elle valait 0.0, ce qui ne fonctionne que parce que \texttt{time.time()} est un grand nombre. Une horloge simulée démarrant à 0 aurait bloqué le premier ordre, et ce bogue n'apparaît qu'avec la nouvelle conception.

D'autres entrées cachées cassent la répétabilité~: les nombres aléatoires sans graine enregistrée~; le parcours d'un \texttt{set} Python, dont l'ordre pour les chaînes change d'un lancement à l'autre~; les threads, dont l'entrelacement modifie l'ordre des événements~; la lecture de la date du jour, d'un fichier susceptible de changer, ou d'un prix issu d'un service en direct. L'ordre des événements de même horodatage a lui aussi besoin d'une règle stable.

Le test de l'ensemble est simple~: lancer deux fois la même entrée et comparer chaque sortie, octet par octet. Il faut en faire un test permanent. En C++, l'injection peut prendre la forme d'un paramètre de modèle, de sorte que le chemin de production n'en paie pas le coût.

\iqwrong{«~Attendre dans le rejeu pour que le temps passe correctement.~» Alors une année de données prend une année, et ce n'est toujours pas déterministe.}

\iqwhy{«~Le même code en recherche et en production~» est un objectif presque partout. Le temps est l'endroit où cela casse généralement en premier.}
\end{solution}

\begin{solution}{iq:iv3:reading-debugging-and-software-design:37}
Générer des ajouts et des annulations aléatoires, garder en parallèle un modèle simple des ordres actifs à côté du carnet, vérifier les invariants après chaque étape, et réduire un échec en supprimant des opérations. Tout tient à deux opérations~: ajouter l'ordre 5 au prix 100, puis annuler l'ordre 5. Le bogue~: \texttt{cancel} laisse un niveau de quantité 0, de sorte que \texttt{best} renvoie un prix que personne ne cote.

Le \Cref{lst:iv3:rd:propcheck} est tout le dispositif, en Python pur avec \texttt{random}. Il n'utilise aucune bibliothèque. Chaque exécution tire 30 opérations sur un petit intervalle d'identifiants et de prix, car les petits intervalles rendent les collisions probables, et les collisions sont là où vivent les bogues. Le modèle est un dictionnaire des ordres actifs. Il est assez simple pour être manifestement juste, et c'est là le point. Après chaque étape, deux faits doivent être vrais~: le meilleur prix égale le prix le plus haut parmi les ordres actifs, et la somme sur les niveaux égale la somme sur les ordres actifs.

\omcode{code/interviews-3/09-reading-debugging-and-software-design/python/rd_prop_check.py}{3}{41}{Un test aléatoire et un réducteur simple.\protect\hfil\protect\break}\label{lst:iv3:rd:propcheck}

La réduction est simple aussi~: essayer de retirer chaque opération à tour de rôle, et garder le retrait si l'échec persiste. Cette réduction ramène 30 opérations aux 2 citées plus haut, que l'on lit d'un coup d'œil. Pour corriger, il faut supprimer le niveau lorsqu'il atteint zéro. Après cela, 1000 exécutions ne trouvent rien. La graine est fixée, de sorte qu'un échec peut toujours être rejoué.

Le générateur ignore les ajouts d'un identifiant actif, donc il ne dit rien de ce cas, et \texttt{Book} le traite mal~: \texttt{add(1, 100, 5)}, \texttt{add(1, 101, 5)} laisse 5 lots bloqués au prix 100. Il faut soit rejeter un identifiant en double, soit laisser le test en générer un et préciser le résultat attendu.

Pourquoi les tests unitaires l'ont-ils manqué~? \texttt{test\_cancel\_lower} annule le plus bas de deux prix, et seul le meilleur est vérifié, de sorte qu'un niveau inférieur vide ne se voit pas. Son auteur a testé ce qu'il avait imaginé. Un test aléatoire n'a pas de telles préférences.

La partie difficile est le choix des invariants, pas la génération des entrées. Les bons invariants sont~: comparer à un modèle lent, vérifier qu'un aller-retour restaure l'état, et vérifier la conservation de la quantité. Une bibliothèque comme Hypothesis\index{Hypothesis} fait mieux la génération et la réduction, et une société qui l'utilise déjà doit s'en servir. En entretien, savoir écrire la version de 30 lignes compte davantage.

\iqwrong{«~Ajouter un test pour l'annulation du meilleur.~» Cela attrape ce bogue, et il est bon de le garder comme test de non-régression. Il ne fait rien pour le suivant, auquel on n'a pas pensé.}

\iqwhy{Les carnets d'ordres ont beaucoup d'états et peu d'entrées simples, c'est le cas idéal pour le test aléatoire. Un candidat qui le propose spontanément se distingue.}
\end{solution}

\begin{solution}{iq:iv3:reading-debugging-and-software-design:38}
Ajouter le champ en option avec une valeur par défaut, et mettre à niveau d'abord tous les lecteurs, puis les écrivains. Retirer l'ancien champ dans l'ordre inverse~: faire cesser aux lecteurs le besoin de ce champ, faire cesser l'envoi aux écrivains, puis marquer son numéro comme jamais réutilisable. Garder les décodeurs des anciens enregistrements aussi longtemps que l'on garde les enregistrements.

L'ordre compte, car on ne peut pas tout basculer d'un coup~: à chaque étape, chaque lecteur doit traiter chaque message qui circule alors sur le bus. Déploiement pour l'ajout~:
\begin{itemize}
\item Modifier le schéma, et donner au nouveau champ une valeur par défaut signifiant «~non connu~».
\item Déployer les six lecteurs. Ils acceptent les anciens et les nouveaux messages, et utilisent la valeur par défaut lorsque le champ est absent.
\item Déployer les écrivains, qui commencent à remplir le champ.
\item Alors seulement, écrire la logique qui repose sur la présence du champ.
\end{itemize}
Pour le retrait, procéder à l'envers~: les lecteurs cessent d'utiliser le champ, puis les écrivains cessent de l'envoyer, puis le schéma marque le champ comme retiré.

Le format doit permettre à un ancien lecteur d'ignorer ce qu'il ne connaît pas, et à un nouveau lecteur de constater ce qui est absent. Un format à étiquettes offre cela directement. La documentation de Protobuf indique que les anciens binaires ignorent simplement un nouveau champ à l'analyse, et que retirer un champ est sûr tant que son numéro n'est pas réutilisé. C'est la règle qui compte~: réutiliser un numéro rend le décodage ambigu, et son guide des bonnes pratiques recommande de réserver les numéros supprimés et de ne jamais ajouter de champ obligatoire.

Les formats à disposition fixe, fréquents sur les chemins sensibles à la latence, demandent plus de soin. Un champ ne peut être ajouté qu'à la fin d'un bloc, car ailleurs il décale la position de chaque champ suivant. L'en-tête du message a alors besoin d'une version, ou d'une longueur pour que les lecteurs puissent sauter ce qu'ils ne comprennent pas. La signification du champ et ses unités font aussi partie du contrat. Passer des cents aux ticks sous le même champ est le pire cas~: chaque décodeur continue de fonctionner et chaque nombre est faux. Il faut plutôt ajouter un nouveau champ pour cela.

Tests à ajouter. Conserver dans le dépôt un jeu de messages enregistrés pour chaque version passée, et les décoder à chaque construction. Faire comparer le schéma avec le précédent par la construction, et la faire échouer en cas de changement incompatible. Pendant le déploiement, faire tourner côte à côte les anciens et les nouveaux décodeurs sur le trafic en direct, et comparer leurs sorties.

Les données enregistrées modifient l'horizon temporel. Un rejeu dans cinq ans doit encore lire les messages d'aujourd'hui. Les schémas ne sont donc jamais supprimés, la version voyage avec les données, et un changement qui paraît sûr pour le trafic en direct peut encore casser la recherche.

\iqwrong{«~Monter de version et déployer tout pendant le week-end.~» Les données enregistrées portent encore l'ancienne version, un retour arrière la réintroduit, et six équipes sont rarement prêtes au même moment.}

\iqwhy{Les sociétés vivent des années avec leurs formats. Connaître l'ordre du déploiement distingue ceux qui l'ont pratiqué de ceux qui n'ont fait que le lire.}
\end{solution}

\begin{solution}{iq:iv3:reading-debugging-and-software-design:39}
Il y a quatre bogues~: un pointeur pendant après un second ordre à un niveau, un effacement dans une boucle for, le mauvais bout du carnet pour les ordres d'achat, et une annulation qui laisse son niveau derrière elle. ASan n'attrape que le premier. Le deuxième est un bogue mémoire qu'il manque, et les deux derniers relèvent d'une logique qu'aucun outil mémoire ne voit.

\textbf{Pointeur pendant.} \texttt{index\_} stocke \texttt{\&level.back()}, un pointeur dans un \texttt{std::vector}, et le \texttt{push\_back} suivant à ce niveau peut réallouer. Séquence~: \texttt{add(100, 1, 5)}, \texttt{add(100, 2, 5)}, \texttt{cancel(1)}. ASan signale \texttt{heap-use-after-free} dans \texttt{cancel}.

\textbf{Effacement dans la boucle.} \texttt{take} efface le niveau de \texttt{levels\_} à l'intérieur de la boucle for par intervalle qui parcourt \texttt{levels\_}, de sorte que l'itérateur de la boucle est invalide lorsqu'il avance. Séquence~: \texttt{add(100, 1, 5)}, \texttt{add(101, 2, 5)}, \texttt{take(7)}. Ici ASan ne signale rien~: le pas se fait dans \texttt{\_Rb\_tree\_increment}, compilé dans la \texttt{libstdc++.so} non instrumentée, et le nœud libéré contient encore de bons pointeurs. Il affiche la bonne exécution. \texttt{-D\_GLIBCXX\_DEBUG} l'arrête, avec un message signalant un itérateur singulier. Un seul outil ne suffit donc jamais.

\textbf{Mauvais côté.} La table ordonne les prix par ordre croissant, donc \texttt{begin()} est le plus bas. Pour les cours vendeurs, c'est le meilleur. Pour les cours acheteurs, c'est le pire. Séquence du côté acheteur~: \texttt{add(100, 1, 5)}, \texttt{add(101, 2, 5)}, \texttt{take(3)}. Elle se fait à 100 et non au cours acheteur de 101. Aucun outil ne le trouve~: seul un test avec un résultat attendu le révèle. Le commentaire de la classe est l'indice.

\textbf{Niveau périmé.} \texttt{cancel} met la quantité à zéro et laisse l'entrée. Après \texttt{add(100, 1, 5)}, \texttt{cancel(1)}, \texttt{best()} renvoie 100 sans rien à ce niveau. \texttt{take} efface chaque niveau épuisé qu'il visite, de sorte que les entrées vides ne s'accumulent que plus profondément dans le carnet, là où aucun \texttt{take} n'atteint, et elles y occupent de la mémoire pour de bon.

Le \Cref{lst:iv3:rd:bookfixed} corrige les quatre. Un \texttt{std::list} par niveau conserve des itérateurs stables. Chaque entrée d'index stocke son prix et son itérateur, de sorte qu'une annulation coûte $O(1)$ plus la recherche du niveau. Le côté est un paramètre de modèle, de sorte qu'un côté acheteur utilise \texttt{std::greater}. \texttt{take} avance avec l'itérateur que renvoie \texttt{erase}. Le test le vérifie sous ASan sur 30 suites aléatoires d'opérations pour chaque côté, contre un modèle Python lent.

\enlargethispage{-3\baselineskip}%
\omcode{code/interviews-3/09-reading-debugging-and-software-design/cpp/book_fixed.hpp}{9}{48}{Le côté corrigé.\protect\hfil\protect\break}\label{lst:iv3:rd:bookfixed}

Sur la conception, je dirais encore quatre choses. Un \texttt{long} pour le prix comme pour la quantité permet de les intervertir sans que le compilateur ne le signale~: il faut donc utiliser des types. Il n'y a aucune vérification des identifiants en double ni des quantités non positives. Une liste alloue à chaque ordre, donc un chemin critique utiliserait plutôt un pool ou une liste intrusive. Enfin, un vrai carnet a besoin d'une interface pour les détails d'une exécution, et pas seulement du volume exécuté.

\iqwrong{«~Il passe sous ASan, donc la gestion de la mémoire est correcte.~» Le deuxième bogue passe sous ASan sur toutes les entrées essayées. Chaque outil ne voit que ce qu'il instrumente.}

\iqwhy{C'est la revue du test à faire chez soi dans sa forme la plus pure. Les recruteurs comptent les bogues, mais regardent aussi si l'on sait montrer chacun par une courte séquence.}
\end{solution}

\begin{solution}{iq:iv3:reading-debugging-and-software-design:40}
Par ordre de gravité~: un état global caché qui mélange les fichiers, une seule ligne fautive qui interrompt toute l'exécution, un plantage sur quantité nulle, des flottants pour de l'argent, et un fichier jamais fermé. La valeur par défaut mutable a l'air mauvaise, mais elle est sans effet ici. Accepter une fonction pure sur des transactions analysées, plus une fine couche autour du fichier.

\textbf{État global.} \texttt{totals} survit entre les appels. Si l'on charge le jour 1 (10 à 100 et 30 à 102, VWAP 101.5), puis le jour 2 (10 à 200), \texttt{vwap(\textquotedbl{}AAA\textquotedbl{})} donne 121.2, en mélangeant les deux jours. Aucune erreur n'apparaît. C'est la pire faute, car le nombre paraît plausible. Cela signifie aussi que la fonction ne peut être testée sans effacer d'abord un état global.

\textbf{Lignes incorrectes non traitées.} Une ligne dont le \texttt{price} vaut \texttt{oops} lève \texttt{ValueError}, et le fichier est perdu à cette ligne. Ainsi, une seule mauvaise ligne fait perdre une journée. Pire, les lignes précédentes sont déjà dans \texttt{totals}, de sorte qu'un appelant qui intercepte l'erreur voit des totaux partiels. Il faut décider d'une politique~: compter et signaler les lignes incorrectes.

\begin{sloppypar}\textbf{Quantité nulle.} Des lignes de \texttt{+10} et \texttt{-10}, comme une correction, provoquent \texttt{ZeroDivisionError} dans \texttt{vwap}.\end{sloppypar}

\textbf{Flottants.} \texttt{sum([0.1] * 10)} ne vaut pas 1.0, donc de petites erreurs s'accumulent sur des millions de lignes, puis \texttt{round} ajoute un arrondi au pair dans le rapport. Les prix arrivent sous forme de texte exact~: il faut donc les garder exacts, sous forme de \texttt{Fraction}, de \texttt{Decimal} ou de ticks entiers.

\textbf{Le fichier n'est jamais fermé.} \texttt{open} n'est pas utilisé avec \texttt{with}. CPython ferme le fichier lorsque l'objet est collecté et lève un \texttt{ResourceWarning}, comme le confirme le test. Rien dans le langage ne garantit quand cela se produit. Le descripteur de fichier est une ressource~: il faut dire quand elle se termine.

\textbf{Valeur par défaut mutable.} \texttt{symbols=[]} n'est jamais modifiée, donc elle ne fait ici aucun mal, et le signaler montre qu'on a lu le code plutôt que reconnu un motif. Demander malgré tout \texttt{None} ou un tuple. Le \texttt{in} sur une liste est aussi linéaire, donc un ensemble serait préférable. Telle qu'elle est écrite, une liste vide signifie «~tout~», une règle surprenante si elle n'est pas documentée.

\textbf{Structure.} \texttt{load} ouvre, analyse, filtre et accumule, et \texttt{report} affiche. Rien ne peut servir isolément.

Le \Cref{lst:iv3:rd:reviewfixed} sépare ces éléments~: \texttt{parse} transforme les lignes en transactions et collecte les erreurs, \texttt{vwaps} est une fonction pure, et \texttt{report} ne gère que le fichier. Les tests couvrent les deux jours donnant chacun leur réponse, les lignes incorrectes comptées avec leur numéro de ligne, un prix écrit en mots, une ligne trop courte ou un prix de \texttt{1/0}, \texttt{None} pour une quantité nulle, dix transactions à 0.1 donnant exactement 0.1 en moyenne, et, sur des transactions aléatoires, le fait que mélanger l'entrée puis appeler une seconde fois ne change rien.

\enlargethispage{-3\baselineskip}%
\omcode{code/interviews-3/09-reading-debugging-and-software-design/python/rd_review_fixed.py}{1}{33}{La version à retenir.\protect\hfil\protect\break}\label{lst:iv3:rd:reviewfixed}

La version corrigée ne garde qu'une paire cumulée par symbole, de sorte que la mémoire reste faible. Elle ne vérifie pas les transactions en double, et en entretien je demanderais ce que garantit le fichier avant de l'ajouter. Un octet NUL arrête encore le lecteur \texttt{csv} lui-même sous Python 3.10~: il faut le traiter comme un fichier corrompu, pas comme une ligne incorrecte.

\iqwrong{«~Commencer par la valeur par défaut mutable.~» C'est la première chose que repère une lecture par motifs, et le moins nuisible ici. Il faut ordonner la revue selon ce qui peut coûter de l'argent.}

\iqwhy{La revue fait une part importante du métier, et une revue de test à faire chez soi évalue le jugement~: ce que l'on met en premier, et la capacité à montrer chaque problème par une entrée.}
\end{solution}
